\section{Latihan}
\label{sec:latihan-bab7}

\begin{enumerate}
    \item \textbf{Konseptual.} Jelaskan hubungan antara
          \eng{mnemonic} dan \eng{opcode}. Mengapa programmer cukup
          menulis \lstasm|MOV AX, BX| tanpa menentukan sendiri byte
          \hex{89 D8}?

    \item \textbf{Konseptual.} Sebutkan tiga sebab yang membuat satu
          \eng{mnemonic} seperti \asm{MOV} memiliki banyak bentuk
          \eng{encoding}. Berikan satu contoh instruksi untuk
          masing-masing sebab.

    \item \textbf{Analisis field.} Sebuah instruksi memiliki byte
          \hex{C7 80 34 12 78 56}. Tentukan:
          \begin{enumerate}
              \item \eng{opcode}-nya dan operasi yang dilakukan;
              \item nilai biner field \code{mod}, \code{reg}, dan
                    \code{r/m};
              \item panjang dan isi \eng{displacement};
              \item panjang dan isi \eng{immediate}.
          \end{enumerate}

    \item \textbf{Encoding manual.} Tuliskan byte kode mesin untuk
          instruksi berikut beserta langkah perhitungannya.
          \begin{enumerate}
              \item \lstasm|MOV DX, 00FFh|
              \item \lstasm|ADD AL, 3|
              \item \lstasm|MOV [BX], AL|
          \end{enumerate}

    \item \textbf{Analisis ModR/M.} Uraikan byte ModR/M berikut
          menjadi ketiga fieldnya, lalu jelaskan operand apa yang
          diwakilinya.
          \begin{enumerate}
              \item \hex{D8}
              \item \hex{44}
              \item \hex{C1}
              \item \hex{16}
          \end{enumerate}

    \item \textbf{Displacement dan immediate.} Jelaskan perbedaan
          peran \eng{displacement} dan \eng{immediate}. Kapan
          keduanya muncul pada sebuah instruksi, dan dapatkah
          keduanya muncul sekaligus? Berikan contoh.

    \item \textbf{Analisis ukuran dan kecepatan.} Dua program memiliki
          hasil akhir sama. Program pertama memakai
          \lstasm|MOV AX, 0001h| dan \lstasm|MOV BX, 0001h|; program
          kedua memakai \lstasm|MOV BX, 0001h| dan \lstasm|MOV AX, BX|.
          \begin{enumerate}
              \item Hitung jumlah byte kedua program.
              \item Sebutkan field yang menjadi penyebab selisihnya.
              \item Kaitkan dengan jumlah byte yang harus diambil
                    prosesor.
          \end{enumerate}

    \item \textbf{Analisis kode dan kesalahan.} Diberikan keluaran
          perintah \code{U} berikut.

\begin{session}[numbers=none]
1A2B:0100 B80500        MOV     AX,0005
1A2B:0103 8908          MOV     [BX+SI],CX
1A2B:0105 CD21          INT     21
\end{session}

          \begin{enumerate}
              \item Jelaskan mengapa baris kedua dapat berakhir pada
                    alamat \hex{103}, bukan \hex{102}.
              \item Mengapa disassembler menampilkan
                    \lstasm|[BX+SI]| dan bukan nama variabel yang
                    mungkin dipakai pada kode sumber?
              \item Sebutkan satu risiko bila hasil \eng{disassembly}
                    langsung dipercaya tanpa memeriksa kode
                    sumbernya.
          \end{enumerate}
\end{enumerate}

\begin{refleksibox}
Jawab pertanyaan berikut pada jurnal praktikum Anda.
\begin{enumerate}
    \item Bagian mana yang paling menantang saat Anda mengisi lembar
          kerja \eng{encoding} secara manual: menentukan
          \eng{opcode}, menyusun ModR/M, atau mengurutkan
          \eng{little-endian}? Jelaskan alasannya.
    \item Setelah melihat hasil perintah \code{U}, apakah dugaan awal
          Anda tentang panjang sebuah instruksi berubah? Bagian mana
          yang paling mengejutkan?
    \item Konsep apa pada bab ini yang masih perlu Anda latih lagi
          sebelum mempelajari mode pengalamatan pada bab berikutnya?
\end{enumerate}
\end{refleksibox}
